%% sections/05-discussion.tex

\section{Discussion}
\label{sec:discussion}

\subsection{What the 12-Point Improvement Measures}

The improvement in mean gate scores from 65/80 in the earliest corpus sessions to
77/80 in the latest reflects two distinct and partially confounded mechanisms. The
first is genuine quality improvement in adventure design and party construction: the
workshop accumulated design knowledge across 49 sessions, and later adventures are
better engineered for the qualities the rubric measures. This is the expected result
of any iterative design system with quality feedback. The second mechanism is rubric
learning: 13 amendments across the corpus progressively expanded the rubric's ability
to recognize quality that existed earlier but could not be scored, particularly in the
Atmospheric Landing and Character Agency dimensions.

These mechanisms cannot be cleanly separated with the current data, but the version-controlled
opening-session analysis (Table~\ref{tab:version-control}) constrains the rubric-learning
hypothesis. If amendments were purely inflating scores for the same underlying quality,
later campaign opening sessions would score higher than earlier ones under the same
rubric version. Instead, C3's opening session scored lower than C1's opening session
even under a later rubric version, confirming that the rubric's evolution is not
uniformly permissive. The amendments introduce new qualifying paths for high-band
scores rather than relaxing existing criteria, and the evidence citation requirement
means that a new path only scores if the session log actually contains evidence of the
described pattern.

\subsection{Cross-Campaign Consistency}

The finding that C5--C7 mean gate scores lie within 0.7 points of the C1--C4 mean is
the paper's strongest evidence for the rubric's archetype independence. These campaigns
differ systematically on every axis of party design: C1--C3 used grief-adjacent parties
with explicit loss-hooks; C5 used a professional-code party with no grief infrastructure;
C6 used a duty-spine party with an institutional deliverable; C7 used a party of
strangers with no designed personal stakes at all. If the rubric's dimensions were
primarily measuring features of grief-adjacent emotional content rather than session
execution quality, C5--C7 would score systematically lower, because they have fewer of
the atmospheric and character moments that the rubric's early anchors were calibrated
against. They do not.

The per-dimension breakdown provides additional granularity. Atmospheric Landing, the
dimension most associated with grief-adjacent emotional design, shows mean scores of
9.4 in binding sessions across all seven campaigns, including C5--C7. The six-path
10-band structure ensures that a session does not require grief-adjacent party design
to achieve high AL scores: the act-without-announcement path, the restraint-in-moment
path, and the arc-convergence path are all achievable in professional-code and
strangers-party contexts, as the corpus demonstrates.

\subsection{Limitations}

\paragraph{No Inter-Rater Reliability.}
Every session in the corpus was scored by a single rater (an AI system following the
rubric's evidence-citation protocol). No inter-rater reliability coefficients
(Cohen's $\kappa$~\cite{cohen1960coefficient}, Krippendorff's $\alpha$~\cite{krippendorff2011computing},
or intraclass correlation~\cite{shrout1979intraclass}) have been established. The
evidence-citation requirement provides a partial safeguard against rater drift (the
score is only as valid as the cited evidence), but it does not guarantee that two
human raters reading the same session log would assign the same band on every
dimension. Establishing inter-rater reliability is the most pressing methodological
priority for future validation.

\paragraph{Single AI System.}
All 49 sessions were run by Claude (Anthropic) as both DM and player-character
simulator. It is possible that the rubric's anchors are calibrated to patterns
specific to this system's generation style --- that another LLM-based DM would produce
sessions with systematically different characteristics that the current anchors
would not capture or would under-score. Multi-system validation (running the same
adventures with different AI DM systems and scoring with the same rubric) is necessary
to test this.

\paragraph{Single Campaign Setting.}
All 49 sessions are set in the Dragonlance campaign setting for D\&D 5e, a system
with DC-based skill checks, hit-point combat, and spell-slot resource management. The
rubric's Mechanical Fairness and Table Readiness dimensions are explicitly designed
around these mechanics. Adapting the rubric to Powered by the Apocalypse games (move-based
resolution, no hit points, no DM-facing encounter balance), storytelling games (no dice,
no resource management), or other TRPG systems would require systematic review of every
dimension's band descriptions and anchor language.

\paragraph{Single Setting for Atmospheric Landing.}
Dragonlance's emotional register (grief, honor, institutional loyalty, loss of memory,
craft-witness) influenced the specific anchor language in Atmospheric Landing,
particularly the paths involving grief-paragraphs and reliquary-attuning. Sessions
in a different setting might produce atmospheric moments that the current six-path
10-band structure does not recognize, requiring additional amendments.

\subsection{Inter-Rater Reliability and Construct Validity}

The corpus has a significant methodological limitation that every reviewer identified:
all 49 sessions were scored by a single rater, and no inter-rater reliability
coefficients have been established for the rubric as a whole. We address this
limitation directly, describe what the existing evaluation structure does and does
not provide as a proxy, and scope the residual gap.

\paragraph{Proposed validation study.}
The concrete next step is a two-independent-evaluator design: two human raters score
the same 15 session logs drawn from across the corpus (sampling at least two sessions
per campaign, at least two from early rubric versions and two from late). Each rater
scores all 8 dimensions independently using the rubric documentation without
coordination. Cohen's $\kappa$~\cite{cohen1960coefficient} is computed per dimension
rather than globally, since the rubric's dimensions differ in band-boundary
specificity and anchor density. The target is $\kappa > 0.6$ per dimension ---
substantial agreement~\cite{landis1977measurement} --- before the rubric is used in a
comparative study with a different AI system or human DMs. Dimensions that fall below
the target would require anchor revision or rater-calibration sessions before
generalization.

\paragraph{Panel agreement as a proxy.}
The corpus's 5-lens panel system --- in which each session is reviewed from five
distinct player perspectives with independent dimension-level scores --- provides
a structural proxy for inter-rater agreement, though it is not equivalent. The five
panel reviewers (Tactician, Roleplayer, Lorekeeper, Improviser, Fresh Face) score
the same session log using dimension-level criteria specific to their lens, producing
five independent weighted scores per session. The spread between the highest-weighted
and lowest-weighted panel scores --- the Tactician and Roleplayer lens scores in most
sessions --- ranged from 0.9 to 1.6 points across the full C1--C4 corpus. This spread
is narrow enough to suggest construct consistency: all five lenses are scoring
something recognizably related to session quality, not five independent and unrelated
constructs.

The panel mean also consistently tracks the gate score: across all 49 sessions,
$r > 0.9$ between the gate score and the panel mean (Pearson correlation, computed
per-session on the shared 0--80 scale). This is evidence of convergent validity for
the rubric's dimensions --- the property-based gate scoring and the lens-weighted
panel scoring converge on the same ranking of session quality even though their
scoring mechanisms are structurally different.

\paragraph{What the proxy does not provide.}
The panel agreement proxy does not answer the inter-rater question because all five
panel reviewers are AI instantiations running the same underlying system. It is
possible that the convergence reflects shared generation biases rather than the
validity of the rubric's anchors. Human evaluators may diverge systematically on
dimensions where the anchor language admits multiple interpretations --- particularly
the 7--9 band of Atmospheric Landing, which requires judgment about whether behavioral
evidence constitutes ``silent reception.'' The proposed two-evaluator study would
surface these divergences directly.

\paragraph{Scope of the limitation.}
We report rubric validity in the context of a single AI-simulation system;
generalization to human DMs or other game systems requires additional validation.
The corpus findings --- the score trajectory, the cross-campaign consistency, the
amendment rate decay --- should be interpreted as describing the rubric's behavior
in this system, not as system-independent measurements of session quality.

\subsection{Implications for TRPG Research Tools}

The Marathon rubric demonstrates that structured, evidence-based, multi-dimensional
scoring of AI-simulated TRPG sessions is feasible at scale and produces stable results
across diverse campaign archetypes. For researchers considering the rubric as a
starting instrument, we offer several practical observations.

The amendment rate provides a useful calibration signal: the corpus produced 13
amendments in 49 sessions, with most amendments concentrated in the Atmospheric Landing
dimension. A new corpus in a different setting or system should expect a similar
amendment cluster in whichever dimension is most sensitive to the new context's
emotional or mechanical idiom. Researchers should plan for approximately 4--6 amendments
in the first 10--15 sessions of a new corpus before the rubric stabilizes.

The forward-only amendment constraint is operationally important. Without it, the
temptation to retroactively rescore sessions to confirm that the rubric's evolution is
``coherent'' would conflate learning with measurement. Accepting that early sessions
are scored under a less refined instrument is necessary to preserve the integrity of
the score series as a longitudinal record.

The gate/panel spread metric proved valuable as a diagnostic for structural artifacts.
When gate and panel scores diverge by more than two points, the source is almost always
a structural mismatch between what the rubric's property-based anchors can capture and
what a particular player-lens can evaluate. The spread is not evidence of a rubric
failure; it is evidence of a difference in what the two instruments are measuring.
